\scanpage{132}
\inserted{\unsure{Enfin|\deleted{Enfin}}} Il faut prouver que si $\lambda\notin K$ i.e. $a^2-4Q\ne0$ (cf lemme plus bas), alors \deleted{\unsure{$\lambda$ est purement}} $Q(\lambda)$ est purement imaginaire, i.e. les $\lambda_\alpha$ sont tous non réels. \deleted{Cela signifie que les} \inserted{Comme $\lambda_\alpha$ est solution} de l'équation $\unsure{T^2|T}-a_\alpha T+Q=0$, cela signifie que $a_\alpha^2-4Q<0$, ce qui résulte du fait que $4Q-a^2$ est tot $\ge0$.

b')$\Rightarrow$b) trivial, \deleted{\unsure{puisque}}

b)$\Rightarrow$c) trivial, car $\mathbb K(\lambda)$ est égal à $\mathbb K$ ou est une ext. quadratique purement imaginaire de $\unsure{\mathbb K|K}$.

c)$\Longrightarrow$a) car si $\lambda_\alpha=\sigma_\alpha(\lambda)$, on aura \deleted{\unsure{que $\overline\lambda_\alpha$ est}} \inserted{$\lambda_\alpha\in\mathbb L$, donc,} $\overline\lambda_\alpha$ est l'unique \deleted{\unsure{$\lambda$}} conjugué de $\lambda_\alpha$ sur $\unsure{\mathbb K|K}$, \deleted{\unsure{et}} donc
\[
\lambda_\alpha\overline\lambda_\alpha=N_{\mathbb L/\mathbb K}(\lambda_\alpha)=N_{\mathbb L/\mathbb K}(\lambda)=\lambda\overline\lambda=Q.
\]

\deleted{\uline{Corollaire} Soit $L\subset\mathbb L$, \inserted{Supposons $Q\in Q^{*+}$, $L$ fini sur $Q$, $L$ non \deleted{\unsure{totalement}} réel (i.e. non tot. réel)} Alors $L$ est engendré par un $\lambda$ tel que $\lambda\overline\lambda=Q$. Si $Q\in Z$, on peut supposer \deleted{\unsure{encore}} que $\lambda$ est entier.}

\uline{Remarque} Soit $K$ une extension finie tot. réelle de $Q$, et $Q\in Q^{\unsure{*+}}$, $Q>1$. Je dis qu'il existe alors une extension quadr. purement imaginaire $L$ de $K$, \deleted{telle que} et un entier $i\ge0$, \deleted{tels que} tels que $\unsure{L|K}$ soit engendré \inserted{sur $Q$} par un élément $\lambda$ tel que $\lambda\overline\lambda=Q^i$ i.e. $|\lambda|^2=Q^i$. Si $Q\in Z$, on peut $\overline{\mathrm m}$ prendre $\lambda$ entier algébrique.

Dém. Soit $a\in K$ tel que $\unsure{K}=Q(a)$. \inserted{on peut supposer $a$ entier.} Soit $i$ assez grand pour que $\deleted{\unsure{Q^i}}4Q^i-a^2$ soit tot. $>0$ (\deleted{\unsure{cela est possible}} certes, car $Q\in Q$, $Q>1$). Soit $\lambda$ solution de l'équation $\lambda^2-a\lambda+Q^i=0$. Il satisfait à la condition voulue; si $Q$ est entier, il est entier, à l'évident.
\mspage{131}
